\documentclass[11pt]{article}
\usepackage[margin=1in]{geometry}
\usepackage{amsmath,amssymb,amsthm,mathtools}
\usepackage{graphicx}
\usepackage[hidelinks]{hyperref}
\usepackage{xcolor}

\newtheorem{theorem}{Theorem}
\newtheorem{lemma}[theorem]{Lemma}
\newtheorem{corollary}[theorem]{Corollary}
\newcommand{\Z}{\mathbb Z}
\newcommand{\T}{\mathbb T}
\newcommand{\e}{\mathrm e}
\newcommand{\spec}{\operatorname{Spec}}

\title{Finite-quotient and eventual-block complementary spectra\\
\large A partial result on Question (Q4) of Gabardo--Lai}
\author{Candidate proof packet for arXiv:1311.2209}
\date{13 August 2026}

\begin{document}
\maketitle

\noindent\textbf{Status.}
Candidate partial result, likely valid pending human review.  The unrestricted
type-II question is not claimed solved.

\section{Source question}

Gabardo and Lai~\cite{GabardoLai} consider integers $N_n\geq2$, put
$P_0=1$ and $P_n=N_1\cdots N_n$, and define
\[
 \eta_n=\frac1{N_n}\sum_{r=0}^{N_n-1}\delta_{r/P_n}.
\]
The associated type-II natural pair is
\[
 \mu=\mathop{*}_{j\geq1}\eta_{2j-1},
 \qquad
 \nu=\mathop{*}_{j\geq1}\eta_{2j},
 \qquad \mu*\nu=\mathcal L_{[0,1]}.
\]
On page 22, Question (Q4) asks whether, for every spectrum $\Lambda_mu$ of
$\mu$, there is a spectrum $\Lambda_\nu$ of $\nu$ such that
$\Lambda_\mu\oplus\Lambda_\nu=\Z$.  The paper notes that the remaining case
is type II.

\begin{center}
 \includegraphics[width=0.96\textwidth]{figures/open_problem_crop.png}
\end{center}

\section{Result and proof intuition}

At finite depth, the Fourier zeros have alternating mixed-radix layers.  Every
difference in a spectrum of the odd head lies in an odd layer; every difference
in a spectrum of the even head lies in an even layer.  Hence two head sums
cannot collide modulo $P_{2k}$.  Their cardinalities multiply to $P_{2k}$, so
the sums give every residue exactly once.  This works for arbitrary finite
spectra, not just the digit spectra displayed in the source paper.

For an infinite spectrum with a genuine head--tail product, the finite head
sum may represent residue $r$ by $r+P_{2k}c_r$ rather than by $r$ itself.
An exact tail tiling absorbs the integer carry $c_r$.  This proves a recursive
reduction and, in particular, Question (Q4) for arbitrary finite spectral
perturbations followed by a canonical tail.

The argument does not pass automatically to a general spectral tree.  The
finite decompositions always converge in the mixed-radix profinite completion,
but the resulting paths need not terminate as ordinary integers; moreover, a
maximal orthogonal tree need not be a complete spectrum~\cite{Dai}.  This is
the remaining obstruction.

\section{Finite complementary spectra}

Define the finite heads
\[
 \mu_k=\eta_1*\eta_3*\cdots*\eta_{2k-1},\qquad
 \nu_k=\eta_2*\eta_4*\cdots*\eta_{2k}.
\]
For a finite measure, a ``spectrum'' below is a frequency set whose
exponentials form an orthonormal basis of its finite-dimensional $L^2$ space.

\begin{lemma}[Fourier-zero layers]\label{lem:layers}
The real zero set of $\widehat\eta_n$ is
\[
 Z(\widehat\eta_n)=P_{n-1}\Z\setminus P_n\Z.
\]
Consequently,
\begin{align*}
 Z(\widehat\mu_k)&=\bigcup_{j=1}^k
   \bigl(P_{2j-2}\Z\setminus P_{2j-1}\Z\bigr),\\
 Z(\widehat\nu_k)&=\bigcup_{j=1}^k
   \bigl(P_{2j-1}\Z\setminus P_{2j}\Z\bigr).
\end{align*}
\end{lemma}

\begin{proof}
The Fourier transform of $\eta_n$ is the geometric sum
\[
 \widehat\eta_n(\xi)=\frac1{N_n}\sum_{r=0}^{N_n-1}
       \e^{-2\pi i r\xi/P_n}.
\]
It vanishes exactly when $\xi/P_n$ is a nonintegral $N_n$-th root frequency,
equivalently when $\xi\in P_{n-1}\Z\setminus P_n\Z$.  The remaining claims
follow by taking products.
\end{proof}

\begin{theorem}[Arbitrary finite spectra tile the quotient]\label{thm:finite}
Let $L_0$ be any spectrum of $\mu_k$ and $C_0$ any spectrum of $\nu_k$.
Choose $\ell_0\in L_0,c_0\in C_0$ and normalize
$L=L_0-\ell_0$, $C=C_0-c_0$.  Then
\[
 L\oplus C=H_k,
\]
where $H_k$ is a complete system of residues modulo $P_{2k}$.  Equivalently,
the addition map
\[
 L\times C\longrightarrow \Z/P_{2k}\Z,qquad
 (\ell,c)\longmapsto \ell+c
\]
is bijective.
\end{theorem}

\begin{proof}
The normalization makes both spectra contain zero.  Orthogonality and
Lemma~\ref{lem:layers} then imply $L,C\subset\Z$.  If $\ell\neq\ell'$ in
$L$, the difference $\ell-\ell'$ belongs to a unique odd layer below $2k$.
Likewise, if $c\neq c'$ in $C$, then $c-c'$ belongs to an even layer below
$2k$.

Suppose
\[
 \ell+c\equiv\ell'+c'\pmod{P_{2k}}.
\]
If both differences are nonzero, then
$\ell-\ell'=(c'-c)+mP_{2k}$ for some $m\in\Z$.  Adding a multiple of
$P_{2k}$ cannot change the first failed divisibility layer below $2k$.
The left side therefore has an odd layer and the right side an even layer, a
contradiction.  If one difference is zero, the other is congruent to zero
modulo $P_{2k}$.  This is also impossible for distinct frequencies: the
displayed zero layers contain no multiple of $P_{2k-1}$ for $\mu_k$ and no
multiple of $P_{2k}$ for $\nu_k$.  Thus the addition map is injective.

Finally,
\[
 |L|=N_1N_3\cdots N_{2k-1},\qquad
 |C|=N_2N_4\cdots N_{2k},
\]
because these are the support cardinalities of the finite measures.  Hence
$|L||C|=P_{2k}$, and injectivity is bijectivity.
\end{proof}

In particular, one may take
\[
 C=-\bigoplus_{j=1}^k
 P_{2j-1}\{0,1,\ldots,N_{2j}-1\},
\]
the negative canonical even-head spectrum from~\cite{GabardoLai}.

\section{A head--tail spectrum lemma}

For $P=P_{2k}$, rescale the tails by $P$:
\[
 \bar\mu^{(k)}=(x\mapsto Px)_\#
    \mathop{*}_{j>k}\eta_{2j-1},
 \qquad
 \bar\nu^{(k)}=(x\mapsto Px)_\#
    \mathop{*}_{j>k}\eta_{2j}.
\]
These are again a natural type-II pair, now for the tail sequence
$N_{2k+1},N_{2k+2},\ldots$.

\begin{lemma}[Product spectra with a finite head]\label{lem:product}
If $L$ is a spectrum of $\mu_k$ and $U$ a spectrum of $\bar\mu^{(k)}$, then
\[
 L\oplus P U
\]
is a spectrum of $\mu$.  Similarly, if $C$ is a spectrum of $\nu_k$ and $V$
a spectrum of $\bar\nu^{(k)}$, then $C\oplus PV$ is a spectrum of $\nu$.
\end{lemma}

\begin{proof}
We prove the odd case.  The digit expansion makes addition from the head and
tail supports one-to-one outside the usual null set of ambiguous endpoint
expansions.  It therefore identifies $L^2(\mu)$ with
$L^2(\mu_k)\otimes L^2(\bar\mu^{(k)})$.

Write a point as $a+b/P$, with $a$ in the head support and $b$ in the
rescaled tail support.  Since $a\in P_{2k-1}^{-1}\Z$ and
$P=P_{2k-1}N_{2k}$, one has $Pa\in\Z$.  Thus
\[
 e_{\ell+Pu}(a+b/P)
   =e_\ell(a)\,e_{u+\ell/P}(b).
\]
For each fixed $\ell\in L$, the set
$\{e_{u+\ell/P}:u\in U\}$ is merely a common frequency translate of the
orthonormal basis $E(U)$ in the tail space.  The displayed functions therefore
form an orthonormal basis of the tensor product: the first factors run through
the basis $E(L)$, and inside each first-factor summand the second factors form
a full tail basis.  The even case is identical.
\end{proof}

\section{The eventual block-product theorem}

\begin{theorem}[Recursive tail reduction]\label{thm:tail}
Fix $k\geq1$ and $P=P_{2k}$.  Let $L,C$ be arbitrary normalized spectra of
$\mu_k,\nu_k$, meaning that $0\in L\cap C$.
Suppose $U,V$ are spectra of the renormalized tails
$\bar\mu^{(k)},\bar\nu^{(k)}$ and
\[
 U\oplus V=\Z.
\]
Then
\[
 \Lambda=L\oplus PU,qquad \Gamma=C\oplus PV
\]
are spectra of $\mu,\nu$, respectively, and
\[
 \Lambda\oplus\Gamma=\Z.
\]
The conclusion is invariant under opposite common translations of
$\Lambda$ and $\Gamma$.
\end{theorem}

\begin{proof}
Spectrality follows from Lemma~\ref{lem:product}.  By
Theorem~\ref{thm:finite}, $H=L\oplus C$ contains exactly one representative
$h_r$ of each residue $r\in\{0,\ldots,P-1\}$.  Write
$h_r=r+Pq_r$ with $q_r\in\Z$.

Given $z=r+Pq\in\Z$, choose the unique head pair
$h_r=\ell+c$.  Since $U\oplus V=\Z$, choose the unique $(u,v)$ satisfying
$u+v=q-q_r$.  Then
\[
 z=(\ell+Pu)+(c+Pv).
\]
Conversely, any two such representations have head sums congruent modulo
$P$, hence the same $r$ and, by the directness in
Theorem~\ref{thm:finite}, the same $\ell,c$.  Tail directness then gives the
same $u,v$.  This proves the exact direct sum.

If the prescribed spectrum is translated by $t$, translate the constructed
complement by $-t$; frequency translation preserves spectrality and leaves
the sum equal to $\Z$.
\end{proof}

\begin{corollary}[Eventually canonical spectra]\label{cor:canonical}
Question (Q4) has an affirmative answer for every prescribed spectrum of the
form
\[
 t+L\oplus P_{2k}U_{\rm can},
\]
where $L$ is an arbitrary normalized finite spectrum of $\mu_k$ and
$U_{\rm can}$ is the canonical odd spectrum of the renormalized tail.  More generally, it holds for
every eventual block-product spectrum whose terminal tail pair already has
complementary spectra.
\end{corollary}

\begin{proof}
Gabardo--Lai prove that the canonical odd tail spectrum and the negative
canonical even tail spectrum have direct sum $\Z$~\cite[Section~3]{GabardoLai}.
Apply Theorem~\ref{thm:tail}, and iterate it for the more general statement.
\end{proof}

\section{Verification and limitations}

The accompanying verifier exhaustively enumerates all normalized odd-head
spectra in the binary depth-two quotient and in the mixed-radix cases
$(2,3,2,2)$, $(2,2,3,2)$, and $(3,2,2,2)$.  In each case it checks direct
tiling by the negative canonical even spectrum.  This is a guard against
layer and modulus mistakes, not part of the proof.

The unrestricted part of (Q4) remains open in this packet.  An arbitrary
spectrum may have a tree structure with unbounded delays before a branch
terminates.  Theorem~\ref{thm:finite} supplies compatible decompositions at
every finite quotient, but compactness only produces a profinite branch.
It does not prove that the branch represents an ordinary integer, nor that a
resulting maximal orthogonal even tree is complete.  Eight focused upgrade
routes, including a greedy complement construction and a proposed analytic
completeness transfer, did not remove this obstruction.

A bounded search through 13 August 2026 used the exact question wording,
source citation graph, and later Cantor/Cantor--Moran spectrum literature.  No
explicit all-spectra resolution of (Q4) was found.  The novelty assessment is
therefore moderate and requires expert confirmation.

\section{Human-review recommendation}

The result is suitable for expert review as a substantial partial theorem.
The most important checks are the zero-layer convention in
Lemma~\ref{lem:layers}, the injectivity modulo $P_{2k}$ in
Theorem~\ref{thm:finite}, and the tensor-basis identification in
Lemma~\ref{lem:product}.  Reviewers should preserve the explicit limitation:
no claim is made for a general nonterminating spectral tree.

\begin{thebibliography}{9}
\bibitem{GabardoLai}
J.-P. Gabardo and C.-K. Lai,
\emph{Spectral Measures Associated with the Factorization of the Lebesgue
Measure on a Set via Convolution}, J. Fourier Anal. Appl. 20 (2014), 453--475;
arXiv:1311.2209.

\bibitem{Dai}
X.-R. Dai,
\emph{Spectra of Cantor measures}, Math. Ann. 366 (2016), 1621--1647;
arXiv:1401.4630.
\end{thebibliography}

\end{document}
